\lysh{18868}{2}
\begin{alist}
\item Είναι γνωστό ότι $\ef (\kappa \cdot 360^\circ + \omega) = \ef \omega$ για κάθε ακέραιο $\kappa$. \\
Επομένως $\ef  500^\circ = \ef (360^\circ + 140^\circ) = \ef  140^\circ$.
\item 
\begin{rlist}
\item Το πρόσημο του τριγωνομετρικού αριθμού $\ef  500^\circ$ είναι το ίδιο με το πρόσημο του τριγωνομετρικού αριθμού $\ef  140^\circ$. \\
Αφού $90^\circ < 140^\circ < 180^\circ$, τότε η τελική πλευρά της γωνίας $140^\circ$ βρίσκεται στο 2ο τεταρτημόριο του τριγωνομετρικού κύκλου. Άρα $\ef  140^\circ < 0$.
\item Είναι λοιπόν $\text{A} = \ef  140^\circ \cdot \hm  250^\circ \cdot \syn  300^\circ$. \\
Από το προηγούμενο ερώτημα έχουμε ότι $\ef  140^\circ < 0$. \\
Αφού $180^\circ < 250^\circ < 270^\circ$, τότε η τελική πλευρά της γωνίας $250^\circ$ βρίσκεται στο 3ο τεταρτημόριο του τριγωνομετρικού κύκλου. Άρα $\hm  250^\circ < 0$. \\
Αφού $270^\circ < 300^\circ < 360^\circ$, τότε η τελική πλευρά της γωνίας $300^\circ$ βρίσκεται στο 4ο τεταρτημόριο του τριγωνομετρικού κύκλου. Άρα $\syn  300^\circ > 0$. \\
Επομένως $\text{A} > 0$.
\end{rlist}
\end{alist}

